% Part: first-order-logic
% Chapter: natural-deduction
% Section: identity

\documentclass[../../../include/open-logic-section]{subfiles}

\begin{document}

\olfileid{fol}{ntd}{ide}

\olsection{带\usetoken{P}{derivation}的\usetoken{S}{identity}}

带等词的推导需要额外的推理规则。

\begin{defish}
\AxiomC{}
\RightLabel{\Intro{\eq}}
\UnaryInfC{$\eq[t][t]$}
\DisplayProof
\hfill
\begin{tabular}{r}
\AxiomC{$\eq[t_1][t_2]$}
\AxiomC{$!A(t_1)$}
\RightLabel{\Elim{\eq}}
\BinaryInfC{$!A(t_2)$}
\DisplayProof
\\[3ex]
\AxiomC{$\eq[t_1][t_2]$}
\AxiomC{$!A(t_2)$}
\RightLabel{\Elim{\eq}}
\BinaryInfC{$!A(t_1)$}
\DisplayProof
\end{tabular}
\end{defish}

在上述规则中，$t$、$t_1$ 和 $t_2$ 都是闭项。\Intro{\eq} 规则允许我们在无任何假设的情况下，直接推导出任何形如 $\eq[t][t]$ 的等式。

\begin{ex}
如果 $s$ 和 $t$ 是闭项，那么 $!A(s), \eq[s][t] \Proves !A(t)$：
\begin{prooftree}
\AxiomC{$\eq[s][t]$}
\AxiomC{$!A(s)$}
\RightLabel{$\Elim{\eq}$}
\BinaryInfC{$!A(t)$}
\end{prooftree}
这通常被称为「等同物的可替换性」原则，或称 Leibniz 律（莱布尼茨律）。
\end{ex}

\begin{prob}
证明 $=$ 既是对称的又是传递的，即给出 $\lforall[x][\lforall[y][(\eq[x][y] \lif
    \eq[y][x])]]$ 和 $\lforall[x][\lforall[y][\lforall[z]((\eq[x][y]
    \land \eq[y][z]) \lif \eq[x][z])]]$ 的推导。
\end{prob}

\begin{ex}
我们从语句
\begin{align*}
& \lexists[x][\lforall[y][(!A(y) \lif \eq[y][x])]]
\intertext{推导语句}
& \lforall[x][\lforall[y][((!A(x) \land !A(y)) \lif \eq[x][y])]]
\end{align*}
我们逆向构造该推导：
\begin{prooftree}
\AxiomC{$\lexists[x][\lforall[y][(!A(y) \lif \eq[y][x])]]
\quad \Discharge{!A(a) \land !A(b)}{1}$}
\DeduceC{$\eq[a][b]$}
\DischargeRule{\Intro{\lif}}{1}
\UnaryInfC{$((!A(a) \land !A(b)) \lif \eq[a][b])$}
\RightLabel{\Intro{\lforall}}
\UnaryInfC{$\lforall[y][((!A(a) \land !A(y)) \lif \eq[a][y])]$}
\RightLabel{\Intro{\lforall}}
\UnaryInfC{$\lforall[x][\lforall[y][((!A(x) \land !A(y)) \lif \eq[x][y])]]$}
\end{prooftree}
现在我们需要用到主假设：既然它是存在公式，我们使用 \Elim{\lexists} 推导出中间结论 $\eq[a][b]$。
\begin{prooftree}
\AxiomC{$\lexists[x][\lforall[y][(!A(y) \lif \eq[y][x])]]$}
\AxiomC{$\Discharge{\lforall[y][(!A(y) \lif \eq[y][c]])}{2}$}
\noLine
\UnaryInfC{$\Discharge{!A(a) \land !A(b)}{1}$}
\DeduceC{$\eq[a][b]$}
\DischargeRule{\Elim{\lexists}}{2}
\BinaryInfC{$\eq[a][b]$}
\DischargeRule{\Intro{\lif}}{1}
\UnaryInfC{$((!A(a) \land !A(b)) \lif \eq[a][b])$}
\RightLabel{\Intro{\lforall}}
\UnaryInfC{$\lforall[y][((!A(a) \land !A(y)) \lif \eq[a][y])]$}
\RightLabel{\Intro{\lforall}}
\UnaryInfC{$\lforall[x][\lforall[y][((!A(x) \land !A(y)) \lif \eq[x][y])]]$}
\end{prooftree}
右上角的子推导利用其假设证明了 $\eq[a][c]$ 和 $\eq[b][c]$。这需要两个独立的推导。$\eq[a][c]$ 的推导如下：
\begin{prooftree}
\AxiomC{$\Discharge{\lforall[y][(!A(y) \lif \eq[y][c]])}{2}$}
\RightLabel{\Elim{\lforall}}
\UnaryInfC{$!A(a) \lif \eq[a][c]$}
\AxiomC{$\Discharge{!A(a) \land !A(b)}{1}$}
\RightLabel{\Elim{\land}}
\UnaryInfC{$!A(a)$}
\RightLabel{\Elim{\lif}}
\BinaryInfC{$\eq[a][c]$}
\end{prooftree}
由 $\eq[a][c]$ 和 $\eq[b][c]$，我们通过 \Elim{\eq} 推导出 $\eq[a][b]$。
\end{ex}

\begin{prob}
给出下列公式的推导：
\begin{enumerate}
\item $\lforall[x][\lforall[y][((\eq[x][y] \land !A(x)) \lif !A(y))]]$
\item $\lexists[x][!A(x)] \land \lforall[y][\lforall[z][((!A(y) \land
    !A(z)) \lif \eq[y][z])]] \lif \lexists[x][(!A(x) \land
  \lforall[y][(!A(y) \lif \eq[y][x])])]$
\end{enumerate}
\end{prob}

\end{document}
